\documentclass[a4paper,10pt]{article}
\usepackage[utf8]{inputenc}
\usepackage[T1]{fontenc}
\usepackage[margin=1in]{geometry}
\usepackage{enumitem}
\usepackage{hyperref}
\usepackage{graphicx}
\usepackage{xcolor}
\usepackage{titlesec}

\definecolor{sectioncolor}{RGB}{44,62,80}
\definecolor{bubblebg}{RGB}{239,239,239}
\definecolor{bubbletext}{RGB}{51,51,51}

\titleformat{\section}
  {\color{sectioncolor}\normalfont\Large\bfseries}
  {}{0pt}{}

\titleformat{\subsection}
  {\color{sectioncolor}\normalfont\large\bfseries}
  {}{0pt}{}

\setlist[itemize,1]{leftmargin=\dimexpr 26pt-.10in}

\usepackage{tikz}
\newcommand{\bubble}[1]{%
    \tikz[baseline=(text.base)]{
        \node[draw=none, fill=bubblebg, rounded corners=8pt, inner xsep=6pt, inner ysep=2pt] (text) {\textcolor{bubbletext}{\strut#1}};
    }%
}

\begin{document}

% Header
\begin{center}
    {\LARGE \textbf{Nathan Fitger}}\\[0.5em]
    \textit{Etudiant en master 2 intelligence artificielle à l'universitée de Bordeaux}\\[1em]
    \href{mailto:nath.fitger@gmail.com}{nath.fitger@gmail.com} \textbar\ 
    \href{https://www.linkedin.com/in/nfitger/}{\includegraphics[height=1em]{ressources/logo-linkedin.png} LinkedIn} \textbar\ 
    \href{https://github.com/0xNatgan}{\includegraphics[height=1em]{ressources/github.png} GitHub}
    \href{tel:+330781613102}{\includegraphics[height=1em]{ressources/appel.png} +33 07 81 61 31 02}
\end{center}

\vspace{1em}

% Photo (top left corner)
\begin{picture}(0,0)
    \put(-20,10){\includegraphics[width=2cm]{ressources/_SDB2250-2.jpg}}
\end{picture}

% Profil
\section*{Profil}
Étudiant en intelligence artificielle passionné par l'informatique et les nouvelles technologies.\@ Je suis à la recherche d'un stage dans le domaine de l'intelligence artificielle à partir de février 2026 pour une durée de 6 mois.\@ Ayant exploré divers domaines tels la robotique, le web et la cybersécurité avant de me spécialiser en IA.\@ J'ai acquis des compétences solides en programmation et en analyse de données, et je suis toujours à la recherche de nouveaux défis pour élargir mes connaissances.

Au fil de mes expériences en tant que responsable de parc informatique, animateur et concepteur de stages, participant à des CTF et ensuite stagiaire en IA, j'ai appris à travailler en équipe et à apprendre de nouveaux outils et langages de programmation afin de m'adapter à la tâche demandée.

\noindent\textbf{Langues:} Anglais C1+ (languaskill) \textbar\ Allemand B2\\
\noindent\textbf{Permis B} 

% Compétences
\section*{Compétences}
\noindent
\begin{tabular}{p{0.95\textwidth}}
\bubble{Python} \bubble{LSP-severs} \bubble{C} \bubble{Machine Learning} \bubble{Deep Learning} \bubble{Java} \bubble{Mkdocs} \bubble{Docker} \bubble{Git} \bubble{ctf} \bubble{IA} \bubble{open source} \bubble{RAG} \bubble{Scikit-learn} \bubble{pytorch} \bubble{Neural Networks} \bubble{Langage Models} \bubble{TCL} \bubble{SQL} \bubble{HTML} \bubble{CSS} \bubble{GLPI} \bubble{Markdown} \bubble{Cybersécurité} \bubble{Robotique} \bubble{Web} \bubble{Linux} \bubble{Windows} \bubble{c++} \bubble{JSON}
\end{tabular}

% Expérience
\section*{Expérience}
\noindent\textbf{Stage IA -- Gertrude (saem)}\\ \hfill \textit{Mai 2025 -- Juillet 2025}
\begin{itemize}
    \item Mise en place d'un outil de documentation de base de code à l'aide de modèles de langage
    \item Conseil sur des outils d'intelligence artificielle
    \item Participation à des projets de recherche en IA lié au traffic urbain
\end{itemize}

\noindent\textbf{Responsable du parc informatique, bénévole, trésorier -- Jeunes sciences Bordeaux}\\ \hfill \textit{2019 -- 2025}
\hspace{0.5em}\begin{itemize}
    \item Gestion et maintenance du parc informatique
    \item Assistance technique aux utilisateurs
    \item Maintenance des équipements informatiques
    \item Gestion des stocks de matériel informatique à l'aide de GLPI
\end{itemize}

\noindent\textbf{Animateur de stages d'Informatique et Robotique -- Jeunes sciences Bordeaux}\\ \hfill \textit{2019 -- 2025}\\
\indent Mise en place d'animations de robotique et d'informatique pour les jeunes durant les vacances scolaires

\noindent\textbf{Job d'été restauration – Bordeaux}\\ \hfill \textit{Été 2021 à 2024}
\begin{itemize}
    \item 2021 Eklo Bordeaux -- Runner service en salle, prise de commande et barman
    \item 2023 Heiko poké Bègles -- Employé polyvalent de restauration
    \item 2024 Island poké -- Employé polyvalent de restauration
\end{itemize}

% Formation
\section*{Formation}
\noindent\textbf{Master intelligence artificielle -- Université de Bordeaux} \hfill \textit{2024 -- 2026}\\
\indent Master intelligence artificielle. Option Base de données avancées

\noindent\textbf{Licence Informatique -- Université de Bordeaux} \hfill \textit{2020 -- 2024}\\
\indent Licence Informatique.

\noindent\textbf{Formation développeur web -- LeWagon} \hfill \textit{sept -- nov 2021}\\
\indent Formation intensive au développement web.

\noindent\textbf{Baccalauréat Scientifique -- Lycée Vaclav Havel} \hfill \textit{2019--2020}\\
\indent Baccalauréat Scientifique.


% Projets
\section*{Projets}
\noindent\textbf{DocGen\_LLM} (\href{https://github.com/0xNatgan/DocGen_LLM}{GitHub})\\
Outil de génération de documentation de code à l'aide de modèles de langage et de RAG.\@

\noindent
\bubble{Python} \bubble{Docker} \bubble{Serveurs de langage} \bubble{LLM} \bubble{RAG} \bubble{documentation} \bubble{mkdocs} \bubble{Markdown} \bubble{HTML} \bubble{open source} \bubble{IA}

\begin{itemize}
    \item Problématique principale: Générer une documentation de code de manière automatisée et agnostique au langage dans le cadre de projet manquants de documentation et/ou ayant une documentation peu claire
    \item DocGen\_LLM est un outil de génération de documentation de code qui utilise des modèles de langage en lien avec des serveurs de langage pour générer une documentation claire et concise de l'ensemble d'une base de code d'un projet.\@ Il est conçu pour être agnostique par rapport aux langages de programmation.\@ L'outil s'intègre facilement dans les outils de lecture de documentation tels que mkdocs.\@ L'outil permet de générer une documentation à ajouter au code ou à utiliser dans une documentation externe utilisant markdown ou HTML
\end{itemize}

\noindent\textbf{Serveur de Language TCL} (\href{https://github.com/0xNatgan/tcl-lsp}{GitHub})\\
Serveur de language TCL pour l'analyse de code.

\noindent
\bubble{TCL} \bubble{Serveur de language} \bubble{TCP} \bubble{stdin/out}

\begin{itemize}
    \item Problématique principale: Dans le cadre d'un besoin d'analyse d'un projet en TCL j'ai développé un serveur de language pour faciliter cette tâche et la rendre possible avec mon projet DocGen\_LLM.\@
    \item Ce serveur de language permet d'analyser le code TCL et d'en extraire des informations selon le standard LSP. ce projet n'est qu'une instance simple d'un serveur de langage TCL\@
\end{itemize}

\noindent\textbf{Quoridor}\\
Jeu de société Quoridor développé en Java avec une interface graphique et des IA.\@ Projet réalisé dans le cadre des cours de l'Université de Bordeaux.

\noindent
\bubble{Java} \bubble{IA} \bubble{Graphique} \bubble{JavaFx} \bubble{Junit}

\begin{itemize}
    \item Problématique principale: Développer un jeu de société Quoridor en Java avec une interface graphique et des IA pour jouer contre l'utilisateur.
    \item Ce projet comprend une implémentation du jeu Quoridor avec une interface graphique en Java et des IA pour jouer contre l'utilisateur ainsi que des statistiques avancées sur les algorithmes d'IA utilisés.\@
\end{itemize}

\end{document}